\chapter{Demonstra\c{c}\~{a}o da Aplica\c{c}\~{a}o} \label{cap:demonstracao}

Este cap\'{i}tulo apresenta os principais fluxos da aplica\c{c}\~{a}o desenvolvida, com o objetivo de demonstrar o funcionamento das funcionalidades implementadas at\'{e} \`{a} fase atual do projeto. A demonstra\c{c}\~{a}o encontra-se organizada de acordo com o percurso natural de utiliza\c{c}\~{a}o da plataforma, desde a autentica\c{c}\~{a}o at\'{e} \`{a} gest\~{a}o e acompanhamento dos processos de averigua\c{c}\~{a}o.

\section{Autentica\c{c}\~{a}o e sele\c{c}\~{a}o de papel} \label{sec:demo-auth}

O primeiro contacto do utilizador com a aplica\c{c}\~{a}o ocorre atrav\'{e}s da p\'{a}gina de autentica\c{c}\~{a}o. Nesta p\'{a}gina, o utilizador introduz as suas credenciais para iniciar sess\~{a}o. Ap\'{o}s uma autentica\c{c}\~{a}o bem sucedida, a aplica\c{c}\~{a}o apresenta os pap\'{e}is dispon\'{i}veis para esse utilizador.

\textbf{TODO: colocar imagem da p\'{a}gina de \textit{login} aqui.}

A sele\c{c}\~{a}o de papel permite que o mesmo utilizador desempenhe diferentes fun\c{c}\~{o}es no sistema, quando aplic\'{a}vel. Esta decis\~{a}o define o contexto da sess\~{a}o e condiciona as funcionalidades dispon\'{i}veis na interface, bem como os endpoints que podem ser acedidos no backend.

\textbf{TODO: colocar imagem da p\'{a}gina de sele\c{c}\~{a}o de papel aqui.}

\section{Listagem e acompanhamento de processos} \label{sec:demo-processos}

Depois da sele\c{c}\~{a}o do papel, o utilizador \'{e} encaminhado para a \'{a}rea principal da aplica\c{c}\~{a}o, cujo conte\'{u}do varia consoante o papel selecionado. Para os pap\'{e}is diretamente envolvidos no acompanhamento de processos, esta \'{a}rea pode apresentar uma vis\~{a}o geral dos processos relevantes, incluindo informa\c{c}\~{a}o resumida como o nome do processo, a localiza\c{c}\~{a}o, a \'{a}rea associada, a data limite e a prioridade.

\textbf{TODO: colocar imagem da dashboard/listagem de processos aqui.}

Quando dispon\'{i}vel, a tabela de processos serve como ponto de entrada para a consulta detalhada de cada processo. Esta abordagem permite que os utilizadores envolvidos no \textit{workflow} acompanhem o estado dos processos de forma centralizada, reduzindo a depend\^{e}ncia de informa\c{c}\~{a}o dispersa por diferentes ficheiros ou canais de comunica\c{c}\~{a}o.

\section{Cria\c{c}\~{a}o de processos} \label{sec:demo-criacao-processos}

A cria\c{c}\~{a}o de processos \'{e} uma funcionalidade associada ao papel de triador. Neste fluxo, o triador preenche a informa\c{c}\~{a}o inicial do processo, incluindo dados como o nome, localiza\c{c}\~{a}o, \'{a}rea, prioridade, data limite, seguradora e tipifica\c{c}\~{a}o. Quando aplic\'{a}vel, o triador pode tamb\'{e}m atribuir manualmente um averiguador e um supervisor ao processo.

\textbf{TODO: colocar imagem do formul\'{a}rio de cria\c{c}\~{a}o de processo aqui.}

Este fluxo representa uma das funcionalidades centrais da aplica\c{c}\~{a}o, uma vez que formaliza a entrada de novos processos no sistema e permite associar, desde o in\'{i}cio, os intervenientes respons\'{a}veis pelo seu acompanhamento.

\section{Detalhe do processo} \label{sec:demo-detalhe-processo}

A p\'{a}gina de detalhe do processo disponibiliza informa\c{c}\~{a}o mais completa sobre um processo espec\'{i}fico. Nesta \'{a}rea s\~{a}o apresentados dados como os intervenientes atribu\'{i}dos, a \'{a}rea, datas relevantes, prioridade, estado atual e informa\c{c}\~{a}o associada ao relat\'{o}rio de averigua\c{c}\~{a}o.

\textbf{TODO: colocar imagem da p\'{a}gina de detalhe do processo aqui.}

Esta p\'{a}gina permite concentrar num \'{u}nico local a informa\c{c}\~{a}o operacional do processo, facilitando o acompanhamento por parte dos diferentes pap\'{e}is envolvidos no \textit{workflow}.

\section{Gest\~{a}o de utilizadores} \label{sec:demo-users}

A aplica\c{c}\~{a}o inclui tamb\'{e}m funcionalidades de administra\c{c}\~{a}o, acess\'{i}veis a utilizadores com o papel de administrador. Atrav\'{e}s desta \'{a}rea, \'{e} poss\'{i}vel consultar utilizadores, criar novos utilizadores, alterar pap\'{e}is e gerir a associa\c{c}\~{a}o de utilizadores a \'{a}reas.

\textbf{TODO: colocar imagem da p\'{a}gina de gest\~{a}o de utilizadores aqui.}

Esta funcionalidade \'{e} relevante para manter o sistema alinhado com a organiza\c{c}\~{a}o interna da CAPT, permitindo refletir altera\c{c}\~{o}es nas equipas, responsabilidades e \'{a}reas de especializa\c{c}\~{a}o.

\section{Perfil e hist\'{o}rico de atividade} \label{sec:demo-perfil-historico}

Por fim, a aplica\c{c}\~{a}o disponibiliza uma \'{a}rea de perfil, onde o utilizador pode consultar informa\c{c}\~{a}o associada \`{a} sua conta e \`{a} sua atividade recente. A exist\^{e}ncia de hist\'{o}rico e registo de intera\c{c}\~{o}es contribui para a rastreabilidade dos processos, permitindo perceber que opera\c{c}\~{o}es foram realizadas ao longo do tempo.

\textbf{TODO: colocar imagem da p\'{a}gina de perfil ou hist\'{o}rico aqui.}

Em conjunto, estes fluxos demonstram a base funcional da plataforma \textit{IPW} e a forma como a aplica\c{c}\~{a}o suporta a gest\~{a}o estruturada dos processos de averigua\c{c}\~{a}o.
